\section{Conclusion}\label{sec:conclusion}

The two exclusions close the stabilizer-code distance intervals at
\((n,k)=(16,5)\) and \((17,3)\). The proofs combine restrictions on a
small supported stabilizer with complete classifications on the
complementary coordinates. Nonadditive quantum codes are outside this
extension criterion, which uses a binary stabilizer subspace.

The remaining entries \(\qc{18,4,6}\) and \(\qc{23,1,9}\) in
Grassl's tables \cite{grassl} provide further targets. Lai's remaining
candidate sets \cite{lai}, Section~VII-D, give concrete weight-enumerator
inputs. Applying the method would require a forced support with a
classifiable complement, followed by kernel and coset-weight tests for
its possible local stabilizers.

The hypotheses of the extension criteria determine which search is
applicable. Theorem~\ref{thm:extension} requires all short stabilizer words to lie in the chosen local
stabilizer. The version for arbitrary degeneracy uses a relative-distance
function that depends on the extension map. The balanced-plane bound
assumes that the entire local stabilizer is the specified plane.

Benchmark data and programs are retained in \pathcode{results/benchmarks/}
and \pathcode{source/benchmarks/}.

\noindent\begin{minipage}{\linewidth}
\paragraph{Acknowledgments.}
We thank an expert reader for comments on an earlier draft.
AI systems assisted with proof exploration, programming, verification, and
writing. The mathematical arguments, exact certificates, classifications,
and search records specify the basis for the conclusions.
\end{minipage}
